\documentclass[10pt, oneside]{article} 
\usepackage{amsmath, amsthm, amssymb, calrsfs, wasysym, verbatim, bbm, color, graphics, geometry}
\usepackage[framemethod=TikZ]{mdframed}
\usepackage{xcolor}

\geometry{tmargin=.75in, bmargin=.75in, lmargin=.75in, rmargin = .75in}  

%%%% Styles %%%%

\newmdenv[
  skipabove=10pt,
  skipbelow=10pt,
  backgroundcolor=blue!3,
  linecolor=blue!60!black,
  linewidth=1pt,
  roundcorner=5pt,
]{thmbox}

\newmdenv[
  skipabove=10pt,
  skipbelow=10pt,
  backgroundcolor=gray!10,
  linecolor=gray!10!black,
  linewidth=1pt,
  roundcorner=5pt,
]{defbox}

\newmdenv[
  skipabove=10pt,
  skipbelow=10pt,
  backgroundcolor=yellow!5,
  linecolor=yellow!60!black,
  linewidth=0.8pt,
  roundcorner=5pt,
]{rembox}

% --- New versions
\newenvironment{theorembox}[1][]{\begin{thmbox}\begin{thm}[#1]}{\end{thm}\end{thmbox}}
\newenvironment{definitionbox}[1][]{\begin{defbox}\begin{defn}[#1]}{\end{defn}\end{defbox}}
\newenvironment{remarkbox}[1][]{\begin{rembox}\begin{rem}[#1]}{\end{rem}\end{rembox}}



\newcommand{\R}{\mathbb{R}}
\newcommand{\C}{\mathbb{C}}
\newcommand{\Z}{\mathbb{Z}}
\newcommand{\N}{\mathbb{N}}
\newcommand{\Q}{\mathbb{Q}}
\newcommand{\Cdot}{\boldsymbol{\cdot}}
\numberwithin{equation}{section}
\newtheorem{thm}{Theorem}[section]
\newtheorem{defn}{Definition}[section]
\newtheorem{conv}{Convention}
\newtheorem{rem}{Remark}
\newtheorem{lem}{Lemma}
\newtheorem{cor}{Corollary}
\newtheorem{proposition}{Proposition}[section]
\newtheorem{example}[thm]{Example}
\usepackage{enumitem}



\title{Stochastic Dynamical Models \\ \large Politecnico di Milano}
\author{Juan Pablo Vallejo Montañez}
\date{Academic Year 2025-2026}

\begin{document}

\maketitle
\tableofcontents

\vspace{.25in}

\section{Introduction}

\begin{definitionbox}[Probability Space]
    A \textbf{probability space} is a measure space \((\Omega,\mathcal{F},P)\), where the measure P satisfies P$(\Omega) = 1$
    \begin{itemize}
    \item $\Omega$: \textbf{Sample space} (set of all possible outcomes or results of that experiment).
    \item $\mathcal{F}$: \textbf{Event space}. Event = subset of outcomes of an experiment to which the probability is assigned.
    \item For $A \in \mathcal{F}$, $P(A)$: \textbf{Probability function}. Assigns to each event in the event space a probability in $[0,1]$.
\end{itemize}
\end{definitionbox}

\noindent
\textit{Example: Coin flip.}

\[
\text{Sample Space: } \{H, T\}, \quad
\text{Random Variable: } X(H)=1,\, X(T)=0, \quad
\text{Probability: } P(\{H\})=P(\{T\})=\frac{1}{2}.
\]

\begin{center}
\texttt{Sample Space} $\longrightarrow$ \texttt{Random Variable} $\longrightarrow$ \texttt{Probability}
\end{center}

\begin{definitionbox}[Sample Space $\Omega$]
$\Omega$ is the set of all possible outcomes of the experiment. 
   
\end{definitionbox}


\begin{definitionbox}[Event Space $\mathcal{F}$]
\(\mathcal{F}\) is a collection of subsets of $\Omega$, called \textbf{events}.It is said that \(\mathcal{F}\) is a \(\sigma\)-algebra on $\Omega$. 

\textbf{Requirements for an event space} (or $\sigma$-algebra) from $\Omega$:
\begin{enumerate}[label=(\roman*)]
    \item Contains both the empty set and $\Omega$: $\emptyset, \Omega \in \mathcal{F}$.
    \item Closed under complement: if $A \in \mathcal{F}$, then $A^c \in \mathcal{F}$.
    \item Closed under (countable) union: if $A_i \in \mathcal{F}$ for $i \in \mathbb{N}$, then \(\displaystyle \bigcup_{i=1}^{\infty} A_i \in \mathcal{F}\).
\end{enumerate}

Conceptually, the event space is the collection of subsets of outcomes of an experiment to which probabilities are asigned.
    
\end{definitionbox}


\begin{definitionbox}[Probability]
The probability measure $P$ is a function $P:\mathcal{F} \rightarrow [0,1]$ that assigns to each event $\mathcal{A} \in \mathcal{F}$ a real number $P(A)$ called this a \textbf{probability}, satisfying;
\begin{enumerate}[label=(\alph*)]
    \item For all $A \in \mathcal{F}$, $0 \le P(A) \le 1$.
    \item $P(\Omega)=1$.
    \item If $A \cap B = \emptyset$, then $P(A \cup B) = P(A) + P(B)$. Generally
    \[
    \text{Give a set of of disjoint events } A_i \cap A_j = \emptyset \text{ for } i \neq j. 
    \]
    \[
    P\!\left( \bigcup_{i=1}^{\infty} E_k \right) = \sum_{i=1}^{\infty} P(E_k).
    \]
\end{enumerate}
    
\end{definitionbox}

\begin{definitionbox}[Random Variable]
A random variable $X$ maps from $\Omega$ to another measurable space $(S, \mathcal{T})$:

\[
X: \Omega \to S.
\]

\[
X^{-1}(S') = \{\omega \in \Omega \mid X(\omega) \in S'\}.
\]

\[
P_X(S') = P(X^{-1}(S')) = P(\{\omega \in \Omega \mid X(\omega) \in S'\}).
\]

If there is some subset $S'$ of $S$ whose probability we want to reason about,  
we measure it by looking at the preimage under $X$ of that subset of $\Omega$.

    
\end{definitionbox}

\textcolor{red}{AÑADIR TODAS LAS IMAGENES Y DEFINICIONES ASOCIADAS A LA VARIABLE ALEATORIA}{}

\section{Discrete-Time Markov Chains}

\begin{definitionbox}[Stochastic Process]
Let $(\Omega,\mathcal{F},\mathbb{P})$ be a probability space and let $I$ be a (countable) state space.
A \textbf{(discrete-time) stochastic process} with values in $I$ is a family
\[
\{X_n\}_{n\in\mathbb{N}_0}, \qquad X_n:\Omega\to I.
\]
For $i\in I$ we write $\mathbb{P}(X_n=i)=\mathbb{P}\{\omega\in\Omega:X_n(\omega)=i\}$.
\end{definitionbox}

\begin{definitionbox}[Markov Chain and Markov Property]
A stochastic process $\{X_n\}_{n\in\mathbb{N}_0}$ with values in $I$ is a
\textbf{Markov chain} if for all $n\in\mathbb{N}_0$ and all $i_0,\dots,i_{n-1},i,j\in I$,
\[
\mathbb{P}(X_{n+1}=j \mid X_n=i,\; X_{n-1}=i_{n-1},\dots,X_0=i_0)
= \mathbb{P}(X_{n+1}=j \mid X_n=i).
\]
We denote the \textbf{one-step transition probabilities} by
\[
p_{ij}(n) := \mathbb{P}(X_{n+1}=j \mid X_n=i), \qquad i,j\in I.
\]
The chain is \textbf{time-homogeneous} if $p_{ij}(n)$ does not depend on $n$; then we write $p_{ij}$.
\end{definitionbox}

\begin{definitionbox}[Transition Matrix and Distributions]
In the time-homogeneous case, the matrix $P=(p_{ij})_{i,j\in I}$ is the
\textbf{transition matrix} and satisfies
\[
p_{ij}\ge 0,\qquad \sum_{j\in I} p_{ij}=1 \quad \text{for every } i\in I
\quad\text{(each row is a probability distribution).}
\]
If $\mu^{(0)}$ is the \textbf{initial distribution} on $I$, $\mu^{(0)}_i=\mathbb{P}(X_0=i)$,
then the law at time $n$ is
\[
\mu^{(n)} = \mu^{(0)} P^{\,n}, \qquad
\mu^{(n)}_j=\mathbb{P}(X_n=j)=\sum_{i\in I}\mu^{(0)}_i\, (P^{\,n})_{ij}.
\]
\end{definitionbox}

\begin{theorembox}[Chapman–Kolmogorov]
For $m,n\in\mathbb{N}$ and $i,j\in I$, let
$p^{(n)}_{ij}:=\mathbb{P}(X_n=j\mid X_0=i)$.
Then
\[
p^{(m+n)}_{ij}=\sum_{k\in I} p^{(m)}_{ik}\, p^{(n)}_{kj},
\qquad\text{equivalently}\qquad P^{\,m+n}=P^{\,m}P^{\,n}.
\]
\end{theorembox}

\begin{remarkbox}[Meaning of a State Transition]
A \emph{state} $i\in I$ represents a possible configuration of the system.
The entry $p_{ij}$ is the probability that, in one time step, the system
moves from state $i$ to state $j$. A row of $P$ lists all possible next states
starting from a fixed current state.
\end{remarkbox}

% -----------------------
% Examples
% -----------------------
\subsection*{Examples}

\paragraph{(1) Fair coin as a Markov chain.}
State space $I=\{H,T\}$. With independent flips,
\[
P=\begin{pmatrix}
\frac{1}{2} & \frac{1}{2}\\[4pt]
\frac{1}{2} & \frac{1}{2}
\end{pmatrix},
\qquad
\mu^{(n)}=\mu^{(0)} P^{\,n}.
\]
Each row is the same because the next outcome does not depend on the current one.

\paragraph{(2) Two–state weather model.}
Let $I=\{\mathrm{S},\mathrm{R}\}$ (sunny/rainy). For $0\le a,b\le 1$,
\[
P=\begin{pmatrix}
1-a & a\\
b   & 1-b
\end{pmatrix}.
\]
Here $a=\mathbb{P}(\text{sun}\to\text{rain})$ and $b=\mathbb{P}(\text{rain}\to\text{sun})$.

\paragraph{(3) Simple random walk on $\mathbb{Z}$.}
State space $I=\mathbb{Z}$. For fixed $p\in[0,1]$ and $q=1-p$,
\[
p_{i,i+1}=p,\qquad p_{i,i-1}=q,\qquad p_{ij}=0 \text{ otherwise.}
\]
Thus $P$ is tridiagonal with constant super/sub-diagonals.

\paragraph{(4) $n$-step transitions.}
For a time-homogeneous chain,
\[
p^{(n)}_{ij}=\mathbb{P}(X_n=j\mid X_0=i)=(P^{\,n})_{ij},\qquad
\sum_{j\in I} p^{(n)}_{ij}=1 \ \ \text{for each } i\in I.
\]







\end{document}
